\documentclass[12pt,a4paper]{article}

\usepackage[utf8]{inputenc}
\usepackage[T1]{fontenc}
\usepackage[english]{babel}
\usepackage{amsmath,amssymb}
\usepackage{geometry}
\usepackage{setspace}
\usepackage{graphicx}
\usepackage{float}
\usepackage{caption}
\usepackage{fancyhdr}
\usepackage{lmodern}
\usepackage{listings}
\usepackage{xcolor}
\usepackage{booktabs}
\usepackage{array}

\geometry{top=2cm,bottom=2cm,left=2.5cm,right=2.5cm}
\setstretch{1.3}

\pagestyle{fancy}
\fancyhf{}
\rfoot{\thepage}
\captionsetup[figure]{position=bottom,justification=centering}

\definecolor{codegray}{rgb}{0.95,0.95,0.95}
\lstset{language=C,basicstyle=\ttfamily\small,backgroundcolor=\color{codegray},
  frame=single,breaklines=true,showstringspaces=false,tabsize=4}

\graphicspath{{images/}}

\newcommand{\img}[3][0.95]{%
  \begin{figure}[H]\centering
    \IfFileExists{images/#2}{\includegraphics[width=#1\textwidth]{#2}}{%
      \fbox{\parbox{0.85\textwidth}{\centering\textbf{Image missing:} \texttt{\detokenize{#2}}}}}%
    \caption{#3}
  \end{figure}}
\newcommand{\code}[1]{\texttt{\small #1}}
\newcommand{\fillme}[1]{\textbf{\textcolor{red}{[#1]}}}

\begin{document}

\begin{titlepage}
\begin{center}
  \IfFileExists{utm_logo_1.jpg}{\includegraphics[width=0.23\textwidth]{utm_logo_1.jpg}\\[0.8em]}{}

  \textbf{MINISTRY OF EDUCATION AND RESEARCH}\\
  \textbf{OF THE REPUBLIC OF MOLDOVA}\\[0.5em]
  \textbf{TECHNICAL UNIVERSITY OF MOLDOVA}\\
  \textbf{Faculty of Computers, Informatics and Microelectronics}\\
  \textbf{Department of Software Engineering and Automatics}\\[2em]

  \vspace{1.5cm}
  \textbf{\Large Tacu Igor \textemdash{} FAF-242}\\[3em]

  \textbf{\Large REPORT}\\[0.5em]
  Laboratory Work No.\ 1\\[0.5em]
  \textbf{Computer Programming}\\[0.5em]
  \textbf{Concurrency in Practice}\\[3em]

  \vfill
  \begin{flushright}
    Lecturer:\\
    Bitca Ernest\\
    DISA, FCIM, UTM
  \end{flushright}
  \vfill

  \large \textbf{Chi\c{s}in\u{a}u \textemdash{} 2026}
\end{center}
\end{titlepage}

\tableofcontents
\newpage

\section{Hardware and Setup}

\begin{table}[H]
  \centering
  \caption{Test machine}
  \renewcommand{\arraystretch}{1.2}
  \begin{tabular}{@{}ll@{}}
    \toprule
    CPU & Apple M1 Pro \\
    Physical cores & 8 (6 performance + 2 efficiency) \\
    Logical processors & 8 (no hyper-threading) \\
    RAM & 16 GB \\
    Disk & internal NVMe SSD (Apple Fabric) \\
    USB drive & 32\,GB flash drive (exFAT), USB; measured about 27\,MB/s (USB~2.0 class) \\
    Operating system & macOS 27.0 (build 26A5425a) \\
    Language / runtime & C (Apple clang 17.0.0, \code{-O2}), POSIX threads \\
    \bottomrule
  \end{tabular}
\end{table}

\noindent\textbf{Data.} \code{numbers.txt}: 50\,000\,000 random integers in
$[-10^6, 10^6]$ (SplitMix64, seed 42), 369\,MB. Reference totals:
negative = 24\,996\,683, positive = 25\,003\,295, zero = 22 (sum 50\,000\,000).
Every run below produced exactly these totals, except the deliberately
unsynchronised version in Task~1b.

\noindent\textbf{Method.} One program (\code{lab1/lab1.c}), task chosen by
command-line arguments. Time = \code{clock\_gettime(CLOCK\_MONOTONIC)} from
thread creation to the last join. Peak memory = \code{ru\_maxrss}.
Tasks 1a/1c: median of 5/3 runs, file in the OS cache.

\section{Task 1: One File}

\subsection{Task 1a: Counting with threads}

The file is cut by byte offset into $T$ parts. A part starts at the first line
start at or after its offset; the next part starts at the same place, so every
line is counted exactly once. Each thread reads only its part with \code{pread}
(1\,MB buffer) and parses it.

\img{task1a.png}{Task 1a: time, speed-up and peak memory vs.\ threads}

\begin{table}[H]
  \centering
  \caption{Task 1a results}
  \begin{tabular}{@{}rrrr@{}}
    \toprule
    Threads & Time (ms) & Speed-up & Peak (MB) \\
    \midrule
    1 & 471 & 1.00 & 2.6 \\
    2 & 242 & 1.94 & 3.7 \\
    4 & 186 & 2.53 & 5.7 \\
    8 & 116 & 4.05 & 9.8 \\
    16 & 104 & 4.51 & 18.0 \\
    32 & 93 & 5.06 & 34.4 \\
    64 & 91 & 5.17 & 48.9 \\
    \bottomrule
  \end{tabular}
\end{table}

\subsection{Task 1b: Three ways to count}

The numbers are loaded into an \code{int} array (not measured); 8 threads (=
logical processors) each take one eighth of the array. The unsynchronised
counters are \code{volatile} so the compiler cannot hide the race in a register.

\begin{table}[H]
  \centering
  \caption{Task 1b: negative / positive / zero totals (expected 24\,996\,683 / 25\,003\,295 / 22)}
  \small
  \begin{tabular}{@{}llrr@{}}
    \toprule
    Version & Run & Time (ms) & Sum of totals \\
    \midrule
    1. shared, no sync & 1 & 63.3 & 12\,574\,364 \\
    & 2 & 30.3 & 11\,647\,577 \\
    & 3 & 32.9 & 12\,113\,996 \\
    2. shared, mutex & median of 3 & 1079.7 & 50\,000\,000 \\
    3. per-thread counters & median of 3 & 11.1 & 50\,000\,000 \\
    \bottomrule
  \end{tabular}
\end{table}

\img[0.55]{task1b.png}{Task 1b: time of the three versions}

\subsection{Task 1c: Producer and consumers}

The main thread is the reader: \code{fgets} line by line, 10\,000 lines per
batch, into a queue (mutex + two condition variables) of at most 100 batches.
Workers pop and parse batches; the reader closes the queue at end of file.
Workers = 1\ldots 64.

\img{task1c.png}{Task 1c: time, speed-up and peak memory vs.\ workers}

\begin{table}[H]
  \centering
  \caption{Task 1c results}
  \begin{tabular}{@{}rrrr@{}}
    \toprule
    Workers & Time (ms) & Speed-up & Peak (MB) \\
    \midrule
    1 & 1765 & 1.00 & 2.7 \\
    2 & 1757 & 1.01 & 2.1 \\
    4 & 1886 & 0.94 & 2.4 \\
    8 & 1761 & 1.00 & 4.5 \\
    16 & 1755 & 1.01 & 2.4 \\
    32 & 1804 & 0.98 & 3.2 \\
    64 & 2269 & 0.78 & 4.2 \\
    \bottomrule
  \end{tabular}
\end{table}

\subsection{Answers to the questions of Task 1}

\begin{enumerate}
  \item \textbf{Where does the speed-up stop growing?} Almost at 8 threads
        (4.05$\times$); 16/32/64 add only 4.5/5.1/5.2$\times$. The CPU has 8
        physical = 8 logical processors, so beyond 8 there is no extra
        hardware to run on. Even at 8 the speed-up is 4, not 8, because 2 of the
        8 cores are slower efficiency cores and the threads share memory bandwidth.
  \item \textbf{Why do 32 and 64 threads not help?} They are not slower
        (93 and 91\,ms vs.\ 116\,ms), only slightly faster: the extra threads
        cannot run in parallel on 8 processors, they only even out the load
        (a thread that lands on a slow core no longer delays the end). Each
        thread adds thread-creation cost and a 1\,MB buffer.
  \item \textbf{Does peak memory depend on the number of threads?} Yes, in
        1a it grows from 2.6 to 48.9\,MB: every thread owns a 1\,MB read buffer
        plus its stack. In 1c it does not (2--4.5\,MB): the data is held in the
        queue, whose size is capped and which here never held more than one batch.
  \item \textbf{Ranking of 1b by time:} 3 per-thread counters (11\,ms)
        $<$ 1 no synchronisation (30--63\,ms) $<$ 2 mutex (1080\,ms).
        Correct: 2 and 3. Version 1 is fast only because it is wrong: the
        threads overwrite each other's increments and about 75\% are lost
        (about 12\,M of 50\,M counted, different every run). Version 2 is
        correct but every element takes a contended lock, so the threads mostly
        wait for each other (100$\times$ slower than version 3). Version 3 has
        no sharing at all; only 8 additions at the end.
  \item \textbf{1c vs.\ 1a.} 1a is faster: 471\,ms (1 thread) to 91\,ms (64)
        against about 1760\,ms for 1c, i.e.\ 3.7--19$\times$ faster. In 1c
        one reader thread does all the file reading with \code{fgets} and
        hands the text over; the reader is the bottleneck, so 1c stops getting
        faster already at \textbf{1 worker} (2 workers: 1.01$\times$; 64
        workers: slower, 2269\,ms, from lock contention).
  \item \textbf{Batches of one line / unbounded queue.} Measured with 8 workers:
        batch = 1 line takes 99.4\,s instead of 1.9\,s (52$\times$ slower)
        because every line costs a \code{malloc}, a lock/unlock and a
        condition-variable signal, more than parsing the line. Batch = 1
        with an unbounded queue: 57\,s, and the queue grew to 71\,624 pending
        items (9\,MB) because the reader outran the workers. With 10\,000-line
        batches the unbounded queue changed nothing (1.9\,s, never more than 1
        batch). In general an unbounded queue lets memory grow without limit
        whenever the producer is faster than the consumers; the bound makes
        the reader wait instead.
\end{enumerate}

\section{Task 2: Many Files on Slow Storage}

The 8 files (6\,250\,000 numbers, about 46\,MB each) are read from the USB
drive; one thread reads a whole file with one \code{read} into one byte buffer,
parses it, and the files are shared round-robin between the threads. Before every run the drive was ejected,
unplugged and plugged back in (Spotlight indexing off). One run per configuration.

\begin{table}[H]
  \centering
  \caption{Task 2: USB drive and internal SSD reference}
  \begin{tabular}{@{}rrrrrrr@{}}
    \toprule
    & \multicolumn{3}{c}{USB drive} & \multicolumn{3}{c}{Internal SSD (reference, \code{F\_NOCACHE})} \\
    \cmidrule(lr){2-4}\cmidrule(lr){5-7}
    Threads & Time (ms) & Speed-up & Peak (MB) & Time (ms) & Speed-up & Peak (MB) \\
    \midrule
    1 & 13686 & 1.00 & 45.7 & 534 & 1.00 & 45.7 \\
    2 & 13066 & 1.05 & 89.8 & 305 & 1.75 & 89.8 \\
    4 & 13226 & 1.03 & 177.4 & 232 & 2.30 & 177.9 \\
    \bottomrule
  \end{tabular}
\end{table}

\img{task2_usb.png}{Task 2 on the USB drive}

\subsection{Answers}
\begin{enumerate}
  \item \textbf{Task 2 on 1 thread vs.\ Task 1a on 1 thread.} Task 1a parses
        the same 50\,M numbers in 471\,ms from the internal disk (data in cache). On the USB
        drive the same work takes 13\,686\,ms (29$\times$ longer; 26$\times$ longer than the 534\,ms from the internal SSD without cache). The parsing
        work is identical, so the extra time is spent waiting for the drive: the
        work is \textbf{I/O-bound} on USB, while Task 1a is CPU-bound.
  \item \textbf{Do more threads help on USB?} No. 2 threads give 1.05$\times$ and 4 threads 1.03$\times$ (13.1\,s and 13.2\,s against 13.7\,s), which is within noise. The USB drive is a single device with a fixed transfer rate
        (about 370\,MB at roughly 27\,MB/s); more threads only
        queue more requests on the same bus. Lecture 1's benefit applies when the
        waiting is on a resource that can serve many requests at once (e.g.\ many
        network servers); here the one shared resource is already saturated by
        one thread, and only the short parsing time can overlap.
\end{enumerate}

\section{Conclusion}

\begin{itemize}
  \item \textbf{CPU-bound:} Task 1a and Task 1b (version 3). Time falls as threads are
        added until the 8 logical processors are used (4$\times$ at 8, then
        flat); on the in-memory array the work scales with cores.
  \item \textbf{Bound by one thread / synchronisation:} Task 1c is limited by
        the single reader (no gain beyond 1 worker); Task 1b version 2 is
        limited by the lock (100$\times$ slower than per-thread counters).
  \item \textbf{I/O-bound:} Task 2 on the USB drive (13.7\,s on one thread, 13.1--13.2\,s on 2 and 4): time is dominated by reading from the drive, and added
        threads do not help.
  \item \textbf{How I know:} CPU-bound work speeds up with threads up to the
        core count and its time is the same whatever the storage; I/O-bound
        work takes much longer than the same parsing from fast storage and does not
        speed up with threads; unsynchronised counters give wrong totals.
\end{itemize}

\section{Answers to the Control Questions}

\begin{enumerate}
  \item \textbf{Physical core vs.\ logical processor.} A physical core is a real
        execution unit; a logical processor is a hardware thread the OS can
        schedule. With hyper-threading one core gives 2 logical processors that
        share its units. This machine has no SMT (8 cores = 8 logical), and the chart
        shows it: the speed-up stops near 8 threads.
  \item \textbf{Boundary between two parts.} Each part starts at the first
        line start at or after its byte offset (if the byte before the offset is
        \code{\textbackslash n}, the offset itself); the end of one part is
        computed by the same function as the start of the next, so they are the
        same position. No line is shared or lost, and all totals sum to 50\,000\,000.
  \item \textbf{How unsynchronised counters lose updates.} \code{counter++} is
        three steps: load, add 1, store. Thread A loads 100; thread B loads 100;
        A stores 101; B stores 101. Two increments, one counted. This happened
        for about 75\% of the increments (12\,M instead of 50\,M).
  \item \textbf{Per-thread counters vs.\ lock.} With a lock only one thread
        increments at a time and the others wait (and the cache line of the
        counter moves between cores). Per-thread counters share nothing, so
        all threads run fully in parallel and are merged once at the end:
        11\,ms vs.\ 1080\,ms.
  \item \textbf{I/O-bound.} The program spends its time waiting for a device, not
        computing. For a USB drive the wait is for its limited transfer speed, a
        single shared resource, so more threads do not make it faster. For a
        server waiting on the network the wait is for remote replies; many
        requests can be in flight at once, so the threads overlap their waits
        and concurrency helps.
  \item \textbf{Equal parts in advance vs.\ a queue.} Equal parts (1a) have no
        queue or locking overhead, and are best when all parts cost the same.
        A queue (1c) balances the load automatically when the items cost
        different amounts or arrive over time (a producer), at the price of
        synchronisation per item (batches reduce it).
\end{enumerate}

\appendix
\section{Build and Run}

\begin{lstlisting}[language=bash]
cd lab1 && make                                  # builds ./lab1
mkdir -p data/parts
./lab1 gen data/numbers.txt                      # 50 000 000 numbers
./lab1 split data/numbers.txt data/parts 8       # numbers_1..8.txt
./lab1 1a data/numbers.txt 8                     # Task 1a, 8 threads
./lab1 1b data/numbers.txt unsync|mutex|local    # Task 1b
./lab1 1c data/numbers.txt 8 10000 100           # Task 1c: workers, batch, queue cap
./run_benchmarks.sh                              # all of Task 1 -> results/
./run_task2.sh /Volumes/USB_NAME                 # Task 2 (asks to replug the drive)
python3 plot.py                                  # charts -> images/
\end{lstlisting}

\end{document}
